\subsection{可解李群}

命题 19. 设 G 为有限维李群。L(G) 为可解的充要条件是 G 具有一个可解开子群。

证明类似于第 5 号命题 12 的证明。

命题 20. 设 G 是 R 或 C 上维数为 n 的单连通可解李群，且 g = L(G)。设 \((\mathfrak{g}_{n}, \mathfrak{g}_{n-1}, \ldots, \mathfrak{g}_{0})\) 是 g 的子代数序列，其维数依次为 n, n-1, \ldots, 0，并满足 \(g_{i-1}\) 是 \(g_{i}\) 的理想（当 i = n, n-1, \ldots, 1 时）。† 设 \(G_{i}\) 是对应于 \(g_{i}\) 的 G 的积分子群。设 \(x_{i}\) 是 \(g_{i}\) 中不属于 \(g_{i-1}\) 的向量。设 \(\phi_{i}\) 是映射

\[
\left(\lambda_ {1}, \lambda_ {2}, \dots , \lambda_ {i}\right) \mapsto \left(\exp \lambda_ {1} x _ {1}\right) \left(\exp \lambda_ {2} x _ {2}\right) \dots \left(\exp \lambda_ {i} x _ {i}\right)
\]

从 \(\mathbf{K}^i\) 到 G。则 \(\phi_n\) 是解析流形的一个同构，且 \(\phi_i(\mathbf{K}^i) = \mathbf{G}_i\) 对所有 \(i\) 成立。

当 \(n = 0\) 时本命题显然成立。我们对 \(n\) 作归纳。设 H 是 G 中满足 \(\mathrm{L(H)} = \mathrm{K}x_n\) 的积分子群。由第 6 节第 6 号命题 14 的推论 1，H 和 \(G_{n-1}\) 都是 G 的单连通李子群，并且作为李群，G 是 H 与 \(G_{n-1}\) 的半直积。因此 \(\lambda \mapsto \exp (\lambda x_n)\) 是从 K 到 H 的同构，而由归纳假设，映射

\[
\left(\lambda_ {1}, \lambda_ {2}, \dots , \lambda_ {n - 1}\right) \mapsto \left(\exp \lambda_ {1} x _ {1}\right) \left(\exp \lambda_ {2} x _ {2}\right) \dots \left(\exp \lambda_ {n - 1} x _ {n - 1}\right)
\]

是从解析流形 \(\mathbf{K}^{n - 1}\) 到解析流形 \(\mathbf{G}_{n - 1}\) 的同构，并将 \(\mathbf{K}^i\times \{0\}\) 映到 \(\mathbf{G}_i\)，其中 \(i = 1,2,\ldots ,n - 1\)。命题得证。

命题 21. 设 G 是 R 或 C 上的有限维单连通可解李群，M 是 G 的一个积分子群。则 M 是 G 的李子群且是单连通的。

继续使用命题 20 中的记号 \(n\)、\(\mathfrak{g}\)、\(\mathfrak{g}_{i}\)、\(x_{i}\)、\(\phi\)，但对 \(x_{i}\) 附加如下条件：设 \(i_{p} > i_{p-1} > \cdots > i_{1}\) 是这些整数 \(i\)，并且它们满足 \(\mathrm{L}(M) \cap \mathfrak{g}_{i} \neq \mathrm{L}(M) \cap \mathfrak{g}_{i-1}\)；于是取

\[
x _ {i _ {k}} \in \mathrm{L} (\mathrm{M}) \cap \mathfrak {g} _ {i _ {k}}
\]

当 \(k = 1, 2, \ldots, p\) 时。对 \(n\) 作归纳，容易看出 \((x_{i_p}, x_{i_{p-1}}, \ldots, x_{i_1})\) 是 \(\mathrm{L}(\mathbf{M})\) 的一个基。设 \(\mathbf{N}\) 是一个单连通李群，使得存在一个同构 \(h\)，它从 \(\mathrm{L}(\mathbf{N})\) 到 \(\mathrm{L}(\mathbf{M})\)。令

\[
y _ {p} = h ^ {- 1} (x _ {i _ {p}}), \dots , y _ {1} = h ^ {- 1} (x _ {i _ {1}}).
\]

由命题 20，映射

\[
\left(\lambda_ {1}, \lambda_ {2}, \dots , \lambda_ {p}\right) \mapsto \left(\exp \lambda_ {1} y _ {1}\right) \left(\exp \lambda_ {2} y _ {2}\right) \dots \left(\exp \lambda_ {p} y _ {p}\right)
\]

是从流形 \(K^{p}\) 到流形 N 的同构。存在从 N 到 G 的李群态射 \(\tau\)，满足 \(h = L(\tau)\) 且 \(\tau(N) = M\)（\(\S 6\)，第 2 号命题 1 的推论 1）。因此 M 是 G 中形如

\[
\begin{array}{r l} \tau ((\exp \lambda_ {1} y _ {1}) \dots (\exp \lambda_ {p} y _ {p})) & = \exp (\lambda_ {1} L (\tau) y _ {1}) \dots \exp (\lambda_ {p} L (\tau) y _ {p}) \\ & = \exp (\lambda_ {1} x _ {i _ {1}}) \dots \exp (\lambda_ {p} x _ {i _ {p}}). \end{array}
\]

因此 \(\mathbf{M} = \phi (\mathbf{T})\)，其中 \(\mathbf{T}\) 是 \(\mathbf{K}^n\) 的一个向量子空间。

命题 22. 设 K = R 或 C。设 V 是有限维向量空间，G 是 GL(V) 的连通可解子群。假设 G 的恒等表示是单的。

\begin{enumerate}
\def\labelenumi{(\roman{enumi})}
\item
  若 \(\mathbf{K} = \mathbf{R}\)，则 \(\dim V \leqslant 2\)，且 \(G\) 是交换的。
\item
  若 K = C，则 dim V = 1。
\item
  设 \(\mathbf{K} = \mathbf{R}\)。则 G 在 \(\mathbf{GL}(V)\) 中的闭包 H 是 \(\mathbf{GL}(V)\) 的连通可解李子群（第 1 号命题 1 的推论 2）。因此 L(H) 是可解的（命题 19）。L(G) 的恒等表示是单的（第 \(\S 6\)，第 5 号命题 13 的推论 2）。所以 dim V \(\leqslant 2\)，且 L(G) 是交换的（第一章第 \(\S 5\)，定理 1 的推论 1 和 4）。因此 G 是交换的。
\item
  设 K = C。令 W 为 V 的非零实向量子空间中在 G 下稳定的集合的一个极小元。由于 G 的恒等表示是单的，由 W 生成的 V 的复向量子空间等于 V。由 (i)，G\textbar W 是交换的。因此 G 是交换的。于是 G 的每个元素都是一个伸缩变换（《代数》第八章第 4 节命题 2 的推论 I），从而 dim V = 1。
\end{enumerate}

推论。设 V 是维数 \textgreater0 的复向量空间，G 是 GL(V) 的连通可解子群。

\begin{enumerate}
\def\labelenumi{(\roman{enumi})}
\item
  存在一个非零元素 \(v\)，属于 \(V\)，使得 \(gv \in \mathbf{C}v\) 对所有 \(g \in G\) 成立。
\item
  存在一个基 \(\mathbf{B}\)，作为 \(\mathbf{V}\) 的基，使得对所有 \(g \in \mathbf{G}\)，\(g\) 关于 \(\mathbf{B}\) 的矩阵是下三角的。
\end{enumerate}

令 \(V_{1}\) 为 V 的非零向量子空间中在 G 下稳定的集合的一个极小元。由命题 22(ii)，\(\dim V_{1} = 1\)。这就证明了 (i)。再对 \(\dim V\) 作归纳，可得存在 V 的在 G 下稳定的向量子空间递增序列 \((\mathrm{V}_1, \mathrm{V}_2, \ldots, \mathrm{V}_n)\)，使得 \(\dim \mathrm{V}_{i+1} / \mathrm{V}_i = 1\) 对所有 \(i < n\) 成立，且 \(\mathrm{V}_n = \mathrm{V}\)；从而得到 (ii)。

